\section{Bandidos Multibrazo, Regret, UCB1 y Thompson Sampling}

\subsection*{a) Regret esperado por exploración aleatoria}

El brazo óptimo es $a^*=2$, con media $\mu^*=0.70$. Las brechas de
suboptimalidad son
\[
    (\Delta_1,\Delta_2,\Delta_3)
    =(0.70-0.20,\,0.70-0.70,\,0.70-0.50)
    =(0.50,\,0,\,0.20).
\]
En $\epsilon$-greedy, la rama exploratoria selecciona uniformemente
entre los tres brazos, incluido el óptimo~\cite[cap.~2]{sutton}.
Condicionado a que el paso sea exploratorio, la brecha media es
$(0.50+0+0.20)/3$. El número esperado de pasos exploratorios es
$T\epsilon=120$. Por linealidad de la esperanza, la contribución
esperada de esos pasos al regret es
\[
    \mathcal{R}_{\mathrm{expl}}(T)
    =T\epsilon\frac{1}{3}\sum_{k=1}^3\Delta_k
    =1000(0.12)\frac{0.70}{3}
    =\boxed{28}.
\]
El regret de la rama de explotación requiere saber cuántas veces la
estimación empírica lleva a elegir un brazo subóptimo.

\subsection*{b) Índices de UCB1}

Con el índice del enunciado y $c=1$,
\begin{align*}
    \mathrm{UCB}_2
    &=0.68+\sqrt{\frac{\ln 50}{30}}
      \approx 0.68+0.361=1.041,\\
    \mathrm{UCB}_3
    &=0.48+\sqrt{\frac{\ln 50}{10}}
      \approx 0.48+0.625=1.105.
\end{align*}
Entre los brazos $2$ y $3$, UCB1 selecciona el brazo $3$ porque la
diferencia en los bonos ($0.264$) supera la desventaja en media
empírica ($0.20$). El bono es mayor para el brazo con menos
observaciones, lo que favorece la exploración bajo
incertidumbre~\cite[cap.~7]{lattimore2020}. Para determinar la
selección entre los tres brazos hace falta el índice del brazo $1$.
El enunciado proporciona su media verdadera, mientras que el cálculo
de UCB1 requiere su media empírica.

\subsection*{c) Posteriores de Thompson Sampling}

\textbf{i) Distribuciones posteriores.}
Tomamos los conteos de este apartado como un escenario distinto del
anterior. Si un brazo con media $\theta_k$ registra $s_k$ éxitos y
$f_k$ fracasos, su verosimilitud es proporcional a
$\theta_k^{s_k}(1-\theta_k)^{f_k}$. Al multiplicar por la densidad
uniforme de $\operatorname{Beta}(1,1)$, la posterior es
$\operatorname{Beta}(1+s_k,1+f_k)$. En particular,
\[
    \theta_2\mid\mathcal{D}\sim\operatorname{Beta}(22,10),
    \qquad
    \theta_3\mid\mathcal{D}\sim\operatorname{Beta}(6,6).
\]

\textbf{ii) Medias, varianzas y exploración.}
Para una posterior $\operatorname{Beta}(\alpha_k,\beta_k)$, la media
es $\alpha_k/(\alpha_k+\beta_k)$ y la varianza es
$\alpha_k\beta_k/[(\alpha_k+\beta_k)^2(\alpha_k+\beta_k+1)]$. Así,
\begin{align*}
    \mathbb{E}[\theta_2\mid\mathcal{D}]&=\frac{22}{32}=0.6875,
    &\operatorname{Var}(\theta_2\mid\mathcal{D})
    &=\frac{22\cdot 10}{32^2\cdot 33}\approx 0.00651,\\
    \mathbb{E}[\theta_3\mid\mathcal{D}]&=\frac{6}{12}=0.50,
    &\operatorname{Var}(\theta_3\mid\mathcal{D})
    &=\frac{6\cdot 6}{12^2\cdot 13}\approx 0.01923.
\end{align*}
El brazo $3$ tiene la posterior más dispersa. Thompson Sampling
extrae un valor de cada posterior y elige el brazo con la mayor
muestra~\cite[cap.~36]{lattimore2020}. Así, el brazo $3$ puede ser
elegido cuando su muestra supera las de los otros brazos, aun
cuando su media posterior es menor que la del brazo $2$.
